\thispagestyle{empty}
\begin{tikzpicture}[remember picture,overlay]
 \fill[indigo] (current page.north west) rectangle ([yshift=-3.6cm]current page.north east);
 \node[anchor=north west,text=white,align=left,inner sep=0pt]
   at ([xshift=2.05cm,yshift=-1.30cm]current page.north west)
   {\fontsize{25}{28}\selectfont\bfseries Contents\\[5pt]
    \fontsize{9.5}{12}\selectfont\color{white!72}\normalfont
    An audit, a route, and one calculation done three ways.};
 \node[anchor=north east,inner sep=0pt]
   at ([xshift=-2.05cm,yshift=-1.25cm]current page.north east)
   {\fontsize{15}{18}\selectfont\bfseries\color{gold}PART X};
\end{tikzpicture}
\vspace*{2.05cm}
{\hypersetup{linkcolor=ink}\renewcommand{\contentsname}{}\setcounter{tocdepth}{2}
\vspace{-3.35cm}\small\setlength{\parskip}{0pt}\tableofcontents}
\vfill
\noindent\begin{tikzpicture}[font=\scriptsize]
 \node[rounded corners=4pt,fill=paper,draw=slate!40,line width=0.7pt,inner sep=10pt,
       text width=16.5cm,align=left] {
   \textbf{\color{indigo}\footnotesize If you want the short version}\\[6pt]
   \textbf{\S1} tells you what is and is not on your shelf. \textbf{\S3} tells you which four pages of
   Axler to read first. \textbf{\S5} is the calculation those four pages let you do. \textbf{\S7} is
   the twelve-week plan.\\[5pt]
   {\color{slate}Two books need buying. Everything else is already in
   \bk{/Users/zaman/Pure\_Maths\_TC\_SelfStudy}.}};
\end{tikzpicture}
\vspace{4pt}
\clearpage
